\documentclass[11pt]{article}
\usepackage[margin=1in]{geometry}
\usepackage{amsmath,amssymb,amsthm,bm}
\usepackage{booktabs}

\newtheorem{proposition}{Proposition}
\newtheorem{remark}{Remark}
\theoremstyle{definition}
\newtheorem{definition}{Definition}

\newcommand{\D}{\mathcal{D}}
\newcommand{\y}{\bm{y}}
\newcommand{\E}{\mathbb{E}}
\newcommand{\R}{\mathbb{R}}
\newcommand{\N}{\mathcal{N}}
\newcommand{\KL}{\mathrm{KL}}
\DeclareMathOperator*{\argmin}{arg\,min}
\DeclareMathOperator*{\argmax}{arg\,max}
\DeclareMathOperator{\softplus}{softplus}
\DeclareMathOperator{\diag}{diag}
\DeclareMathOperator{\Cov}{Cov}
\DeclareMathOperator{\Var}{Var}
\DeclareMathOperator{\tr}{tr}

\title{The Inference Trained in \texttt{pulse\_regressor.py}:\\
Posterior of the Cascade Profile $(\ln a, \ln b)$ from DOM Pulses\\
with Known Vertex, Energy and Shower Axis}
\author{}
\date{2026-09-30}

\begin{document}
\maketitle

\section{Setting and data}

\paragraph{Simulation.}
GENIE $\nu_e$ final states, charged current (CC) and neutral current (NC), are propagated with Geant4 through
homogeneous ice by OMSim. Each event is placed at a vertex $\bm v\in\R^3$ drawn uniformly in $[-8,8]^3$\,m and
rotated by a rotation drawn from the Haar measure on $SO(3)$; the interaction happens at $t_v=0$. The detector is
a $3\times3\times3$ grid of 27 IceCube DOMs (one PMT each) with 7\,m spacing, $\bm r_d\in\{-7,0,7\}^3$\,m,
$d=1,\dots,27$; each records the arrival time of every detected photon. The sample used here has 10\,297 events
(5\,981 CC, 4\,316 NC), of which a fixed random 20\,\% (2\,059) is held out for validation. The classes differ
strongly in visible energy (median 15.9\,GeV for CC, 4.6\,GeV for NC) and in light (median 380 and 98 hits).

\paragraph{Observation.}
Module $d$ records $n_d\in\mathbb{N}_0$ hits at times $t_{d1}\le\dots\le t_{dn_d}$. With the bin edges
$\tau_k = 4\,\text{ns}\,(k-1)$ for $k=1,\dots,K$, $K=48$, and $\tau_{K+1}=+\infty$, its time profile and
summaries are
\begin{align}
  \bm n_d &= (n_{d1},\dots,n_{dK}),\qquad n_{dk} = \#\{j : \tau_k \le t_{dj} < \tau_{k+1}\},
  \label{eq:hist}\\
  \bm q_d &= \big(t_{d1},\ t^{(0.1)}_d,\ t^{(0.5)}_d,\ t^{(0.9)}_d,\ \sigma_d\big),
\end{align}
where $t^{(p)}_d$ is the empirical $p$-quantile and $\sigma_d$ the standard deviation of $\{t_{dj}\}_{j=1}^{n_d}$
(times before $\tau_1=0$ do not occur; the last bin collects all times $\ge\tau_K$). The data of one event are the
set of hit modules,
\begin{equation}
  \D = \big\{(\bm r_d,\ n_d,\ \bm q_d,\ \bm n_d)\ :\ n_d>0\big\},\qquad N = \sum_{d} n_d .
  \label{eq:data}
\end{equation}

\paragraph{Labels.}
Cherenkov light is emitted along the charged tracks, so the profile the light samples is the charged track length
per unit depth. From the Geant4 step points of every charged secondary (e$^\pm$, $\mu^\pm$, $\pi^\pm$, p,
K$^\pm$, $\alpha$) we form segments $i=1,\dots,n_{\mathrm{seg}}$ with midpoints $\bm r_i$ and lengths $w_i>0$, and
write $\langle q\rangle = \sum_i w_i q_i/\sum_i w_i$ for the length-weighted mean of any per-segment quantity
$q_i$. The shower axis is the unit vector maximising the length-weighted second moment about the vertex,
\begin{equation}
  \hat{\bm u} = \argmax_{\bm u\in S^2}\ \sum_i w_i\,\big((\bm r_i-\bm v)\cdot\bm u\big)^2
  = \text{top eigenvector of } M=\sum_i w_i(\bm r_i-\bm v)(\bm r_i-\bm v)^\top,
  \label{eq:axis}
\end{equation}
with its sign fixed by $\sum_i w_i(\bm r_i-\bm v)\cdot\hat{\bm u}>0$. The depth of segment $i$ is
$s_i = (\bm r_i-\bm v)\cdot\hat{\bm u}$; segments with $s_i\le0$ are discarded. The Gamma density
\begin{equation}
  g(s;a,b) = \frac{s^{a-1}e^{-s/b}}{\Gamma(a)\,b^{a}},\qquad s>0,\ a,b>0,\qquad \int_0^\infty g\,ds = 1,
\end{equation}
is fitted to the length-weighted depths by maximum likelihood,
\begin{equation}
  (a_e,b_e) = \argmax_{a,b>0}\ \sum_{i:\,s_i>0} w_i\Big[(a-1)\ln s_i - \frac{s_i}{b} - \ln\Gamma(a) - a\ln b\Big],
  \label{eq:label}
\end{equation}
whose stationarity conditions reduce to one scalar root and a closed form,
\begin{equation}
  \ln a - \psi(a) = \ln\langle s\rangle - \langle\ln s\rangle,\qquad b = \langle s\rangle/a,
  \label{eq:mle}
\end{equation}
with $\psi=\Gamma'/\Gamma$ the digamma function; the right-hand side is positive by Jensen's inequality and
$a\mapsto\ln a-\psi(a)$ decreases strictly from $+\infty$ to $0$, so the root is unique. The target of event $e$ is
\begin{equation}
  \y_e = (\ln a_e,\ \ln b_e)\in\R^2 .
\end{equation}
Events with fewer than five charged tracks have no label and are dropped. $a$ sets the shape (relative width
$1/\sqrt a$, mode $(a-1)b$), $b$ the metre scale; the mean depth is $ab$ and the light constrains it much better
than $a$ and $b$ separately, which is the $a$--$b$ degeneracy visible in every result below. The same track fit
provides the axis $\hat{\bm u}_e$ that is handed to the network as known.

\paragraph{Standardised target.}
Let $\bm m=(0.776,-0.238)$ and $\bm s=(0.450,0.435)$ be the componentwise mean and standard deviation of the labels
over the training split. The network is trained on $\bm z = \diag(\bm s)^{-1}(\y-\bm m)$, an affine and invertible
change of variables made only so that a freshly initialised head starts at the prior (mean 0, width 1 in both
components); Section~6 maps back. Everything below is stated for $\y$.

\begin{remark}[Label noise]
Equation~\eqref{eq:label} is a finite-sample fit, $\y_e = \y_e^{\mathrm{true}}+\bm\epsilon_e$, and a shower of a
few GeV has few tracks, so $\bm\epsilon_e$ is not small. The network learns the posterior of the \emph{labelled}
$\y$, whose spread includes $\Cov(\bm\epsilon)$.
\end{remark}

\section{The inference problem}

\begin{definition}[Known, unknown, marginalised]
\leavevmode
\begin{itemize}
  \item Known per event: $\bm\kappa = (\bm v,\ t_v=0,\ E,\ \hat{\bm u})$, the vertex, the visible energy
        $E=E_{\mathrm{vis}}$ (the summed kinetic energy of the final state) and the shower axis, all taken from the
        simulation truth. In a reconstruction they would come from the vertex, energy and direction fits.
  \item Target: $\y = (\ln a,\ln b)\in\R^2$.
  \item Marginalised: everything else that shaped the event, in particular the class (CC or NC) and the
        shower-to-shower fluctuations of the Geant4 history.
\end{itemize}
\end{definition}

\paragraph{Data-generating process.}
The simulation draws a shower history $\omega\sim p_{\mathrm{gen}}$ (final state, energy, vertex, rotation, every
Geant4 step) and then the detector response $\D\sim p(\D\mid\omega)$. With $T(\omega)=\y$ the label map of
Section~1 and $K(\omega)=\bm\kappa$ the known quantities, the joint law is the push-forward
\begin{equation}
  p(\y,\bm\kappa,\D) = \int p(\D\mid\omega)\,\delta\big(\y-T(\omega)\big)\,\delta\big(\bm\kappa-K(\omega)\big)\,
  p_{\mathrm{gen}}(\omega)\,d\omega .
  \label{eq:joint}
\end{equation}
Throughout, $\bm Y$ denotes the random target with this law and $\y$ one of its values.

\paragraph{Goal.}
For an observed event $(\D^\ast,\bm\kappa^\ast)$ with $p(\D^\ast,\bm\kappa^\ast)>0$, the posterior of the profile
\begin{equation}
  p(\y\mid\D^\ast,\bm\kappa^\ast) = \frac{p(\y,\bm\kappa^\ast,\D^\ast)}{p(\D^\ast,\bm\kappa^\ast)} .
  \label{eq:post}
\end{equation}
It is intractable (it marginalises every history compatible with $\y$), so it is approximated by an amortised
parametric family $q_\phi(\y\mid\D,\bm\kappa)$: one parameter vector $\phi$ serves every event.

\begin{remark}[Mixture prior]
The class is never an input. The prior implied by the training sample is the class mixture
$p(\y\mid\bm\kappa) = \sum_c\pi_c\,p(\y\mid c,\bm\kappa)$ with $\pi_{\mathrm{CC}}:\pi_{\mathrm{NC}} = 58:42$, so where
the light says little the posterior mean is pulled toward the mixture mean, not toward the event's own class. This
matters for the NC/CC use of the posterior (Section~7).
\end{remark}

\section{Objective}

\paragraph{From KL to negative log-likelihood.}
Choose $\phi$ to minimise the expected Kullback--Leibler divergence from the true posterior to the model. With
$\KL(p\,\|\,q)=\int p\ln p-\int p\ln q$ and the conditional entropy
$H(\bm Y\mid\D,\bm\kappa) = -\E[\ln p(\y\mid\D,\bm\kappa)]$,
\begin{align}
  \mathcal R(\phi) &= \E_{p(\D,\bm\kappa)}\Big[\KL\big(p(\cdot\mid\D,\bm\kappa)\,\big\|\,q_\phi(\cdot\mid\D,\bm\kappa)\big)\Big] \notag\\
  &= \underbrace{-\E_{p(\y,\bm\kappa,\D)}\big[\ln q_\phi(\y\mid\D,\bm\kappa)\big]}_{\text{expected conditional NLL}}
    \;-\; \underbrace{H(\bm Y\mid\D,\bm\kappa)}_{\text{independent of }\phi}.
  \label{eq:kl}
\end{align}
Only samples $(\y_e,\bm\kappa_e,\D_e)$ are needed to estimate the first term; the density $p(\y\mid\D,\bm\kappa)$
itself never appears.

\paragraph{Gaussian family.}
The network outputs five numbers per event, $(\mu_1,\mu_2,\ell_{11},\ell_{21},\ell_{22})=\bm o_\phi(\D,\bm\kappa)$, and defines
\begin{align}
  \bm\mu &= (\mu_1,\mu_2)^\top,\qquad
  L = \begin{pmatrix} L_{11} & 0\\ L_{21} & L_{22}\end{pmatrix}, \label{eq:head}\\
  L_{11} &= \softplus(\ell_{11})+\varepsilon,\qquad L_{21}=\ell_{21},\qquad L_{22}=\softplus(\ell_{22})+\varepsilon, \notag
\end{align}
with $\softplus(x)=\ln(1+e^x)$ and $\varepsilon=10^{-3}$, so that $L$ has a positive diagonal and
$\Sigma=LL^\top$ is symmetric positive definite for every output. Then
$q_\phi(\y\mid\D,\bm\kappa)=\N\big(\y;\bm\mu_\phi,\Sigma_\phi\big)$. Because $\ln\det\Sigma = 2\ln L_{11}+2\ln L_{22}$
and $(\y-\bm\mu)^\top\Sigma^{-1}(\y-\bm\mu)=\|\bm z\|^2$ for $\bm z=L^{-1}(\y-\bm\mu)$, the per-event loss is the
exact bivariate Gaussian negative log-likelihood, evaluated by one forward substitution:
\begin{align}
  -\ln q_\phi(\y\mid\D,\bm\kappa) &= \tfrac12\|\bm z\|^2 + \ln L_{11} + \ln L_{22} + \ln 2\pi, \label{eq:nll}\\
  z_1 &= \frac{y_1-\mu_1}{L_{11}},\qquad z_2 = \frac{(y_2-\mu_2) - L_{21}z_1}{L_{22}} . \notag
\end{align}

\begin{proposition}[What the Gaussian head learns]
Fix $(\D,\bm\kappa)$ and suppose $\bm Y\mid\D,\bm\kappa$ has finite second moments, with mean
$\bm m_\D=\E[\bm Y\mid\D,\bm\kappa]$ and covariance $C_\D=\Cov[\bm Y\mid\D,\bm\kappa]\succ0$. Over all
$\bm\mu\in\R^2$ and symmetric positive definite $\Sigma$,
\begin{equation}
  \argmin_{(\bm\mu,\Sigma)}\ \E\big[-\ln\N(\bm Y;\bm\mu,\Sigma)\,\big|\,\D,\bm\kappa\big] = (\bm m_\D,\ C_\D),
\end{equation}
and the minimiser is unique, whether or not the posterior is Gaussian.
\end{proposition}
\begin{proof}
Since $\E[(\bm Y-\bm\mu)(\bm Y-\bm\mu)^\top\mid\D,\bm\kappa] = C_\D+(\bm m_\D-\bm\mu)(\bm m_\D-\bm\mu)^\top$ and
$(\y-\bm\mu)^\top\Sigma^{-1}(\y-\bm\mu)=\tr\big(\Sigma^{-1}(\y-\bm\mu)(\y-\bm\mu)^\top\big)$,
\begin{align}
  \E\big[-\ln\N(\bm Y;\bm\mu,\Sigma)\,\big|\,\D,\bm\kappa\big] &= \tfrac12\tr(\Sigma^{-1}C_\D)
  + \tfrac12(\bm m_\D-\bm\mu)^\top\Sigma^{-1}(\bm m_\D-\bm\mu) \notag\\
  &\quad + \tfrac12\ln\det\Sigma + \ln 2\pi .
\end{align}
For every $\Sigma\succ0$ the second term is non-negative and vanishes only at $\bm\mu=\bm m_\D$. With
$P=\Sigma^{-1}$ the remaining terms are $f(P)=\tfrac12\tr(PC_\D)-\tfrac12\ln\det P+\ln2\pi$, strictly convex on the
positive definite cone because $-\ln\det$ is; its gradient $\nabla_P f=\tfrac12(C_\D-P^{-1})$ vanishes only at
$P=C_\D^{-1}$, i.e.\ $\Sigma=C_\D$.
\end{proof}
Hence, with enough capacity and data, $\bm\mu_{\hat\phi}\approx\E[\bm Y\mid\D,\bm\kappa]$ and
$\Sigma_{\hat\phi}\approx\Cov[\bm Y\mid\D,\bm\kappa]$: one forward pass returns the posterior mean and covariance
of $(\ln a,\ln b)$ given the pulses, the vertex, the energy and the axis. The same holds in the standardised
variables, because the affine map changes the loss only by the constant $\ln s_1+\ln s_2$.

\paragraph{Empirical objective and regularisation.}
With $M$ training events, independent draws from~\eqref{eq:joint},
\begin{equation}
  \hat\phi = \argmin_\phi\ \frac1M\sum_{e=1}^{M}\big[-\ln q_\phi(\y_e\mid\D_e,\bm\kappa_e)\big] + \frac\lambda2\|\phi\|_2^2 ,
  \label{eq:emp}
\end{equation}
the Monte-Carlo estimate of~\eqref{eq:kl} plus a Gaussian prior on $\phi$. It is optimised with AdamW (decoupled
weight decay $10^{-2}$, an approximation of the $L_2$ term), a one-cycle learning-rate schedule with peak $10^{-3}$,
gradient-norm clipping at 1 and minibatches of 64 events. The training set is 80\,\% of the events; the epoch with
the lowest validation NLL is kept and is the model evaluated.

\section{The network, as implemented}

\paragraph{Tokens.}
Each hit module $d$ is one token with a space-time position and a feature vector. With $c$ the vacuum speed of
light and $n_g=1.35$ the group index of ice,
\begin{align}
  \bm\xi_d &= \big(\bm r_d-\bm v\ [\mathrm{km}],\ t_{d1}\ [\mu\mathrm s]\big)\in\R^4, \label{eq:xi}\\
  \bm f_d &= \Big(\ln(1+n_d),\ \frac{\bm q_d}{100\,\mathrm{ns}},\ \frac{t^{\mathrm{res}}_d}{100\,\mathrm{ns}},\
             \frac{|\bm r_d-\bm v|}{10\,\mathrm m},\ \frac{\bm n_d}{n_d},\ \ln E,\ \cos\psi_d\Big)\in\R^{10+K}, \label{eq:f}\\
  t^{\mathrm{res}}_d &= t_{d1} - \frac{n_g}{c}\,|\bm r_d-\bm v| ,\qquad
  \cos\psi_d = \frac{(\bm r_d-\bm v)\cdot\hat{\bm u}}{|\bm r_d-\bm v|} .
\end{align}
The known quantities enter as follows. The vertex fixes the frame: every position is $\bm r_d-\bm v$, so
translation invariance is built in, and the residual time $t^{\mathrm{res}}_d$ is the delay of the first photon
with respect to unscattered light from the vertex. The energy enters as $\ln E$ on every token (the same number on
all tokens of an event). The axis enters as $\cos\psi_d$, the cosine of the angle between the module direction and
the shower axis, which places every module relative to the Cherenkov cone; it is the only piece of
$\hat{\bm u}$ the model needs, since a shower is symmetric about its axis. With $K=48$ the feature vector has
$58$ entries. The normalised histogram $\bm n_d/n_d$ carries the pulse shape, $\ln(1+n_d)$ its size.

\paragraph{Encoder (Neptune).}
Let $n=|\D|$ be the number of hit modules and $D_{\mathrm{tok}}=64$.
\begin{enumerate}
  \item Per-token embedding $\bm x^{(0)}_d = \mathrm{MLP}_1(\bm f_d)\in\R^{D_{\mathrm{tok}}}$. Because $n\le27$ is
        below the tokenizer's capacity, no farthest-point pooling takes place: every hit module is its own token.
  \item Three pre-norm transformer layers, $X\leftarrow X+\mathrm{Attn}(\mathrm{N}(X))$,
        $X\leftarrow X+\mathrm{SwiGLU}(\mathrm{N}(X))$, with RMS normalisation $\mathrm N$ and four heads of
        dimension $d_h=16$. For one head with queries $\bm q_d$, keys $\bm k_{d'}$ and values $\bm v_{d'}$, the
        four-dimensional rotary encoding applies an orthogonal block rotation $R(\bm\xi)$ with
        $R(\bm\xi)^\top R(\bm\xi')=R(\bm\xi'-\bm\xi)$, so that
        \begin{align}
          A_{dd'} &= \frac{\big\langle R(\bm\xi_d)\bm q_d,\ R(\bm\xi_{d'})\bm k_{d'}\big\rangle}{\sqrt{d_h}}
                  = \frac{\bm q_d^\top R(\bm\xi_{d'}-\bm\xi_d)\,\bm k_{d'}}{\sqrt{d_h}},\\
          \mathrm{Attn}(X)_d &= \sum_{d'}\operatorname{softmax}_{d'}(A_{dd'})\,\bm v_{d'} ,
        \end{align}
        so the scores depend on positions only through relative space-time offsets. An additive Fourier encoding
        of the absolute (vertex-relative) position is also present. Both encodings were designed for km-scale
        detectors; their frequencies are multiplied by 30 so that they resolve the metres of this 30\,m geometry.
  \item Attention pooling with a learned query $\bm a$,
        \begin{equation}
          \beta_d = \operatorname{softmax}_{d}\Big(\frac{\bm a^\top W_k\bm x^{(3)}_d}{\sqrt{D_{\mathrm{tok}}}}\Big),\qquad
          \bm h = W_o\sum_{d}\beta_d\,W_v\bm x^{(3)}_d ,
        \end{equation}
        followed by the two event-level scalars the pooling would otherwise normalise away,
        $\bm c = \bm h \oplus \big(\ln(1+n),\ \ln(1+N)\big)$.
  \item Head $\bm o_\phi = \mathrm{MLP}_2(\bm c)\in\R^5$, mapped to $(\bm\mu,L)$ by~\eqref{eq:head}.
\end{enumerate}
In total $1.8\times10^5$ parameters. The map $(\D,\bm\kappa)\mapsto(\bm\mu,\Sigma)$ is a function of the
\emph{set} of hit modules: invariant to their order and defined for any $n\ge1$.

\section{Rotation augmentation}

Let $R\in SO(3)$ and define the rotated event
\begin{equation}
  R\D = \big\{(\bm v+R(\bm r_d-\bm v),\ n_d,\ \bm q_d,\ \bm n_d)\ :\ n_d>0\big\},\qquad
  R\bm\kappa = (\bm v,\ t_v,\ E,\ R\hat{\bm u}),
\end{equation}
i.e.\ module positions and axis rotated about the vertex, times, pulses and energy unchanged. The ice is isotropic,
so the response of the detector geometry $G=\{\bm r_d\}$ to a shower history $\omega$ and the response of the rotated
geometry $RG$ to the rotated history $R\omega$ have the same law up to relabelling of positions,
$p_{RG}(R\D\mid R\omega)=p_G(\D\mid\omega)$; the label is invariant, $T(R\omega)=T(\omega)$; and $p_{\mathrm{gen}}$
is rotation invariant by construction. Therefore, if $(\y,\bm\kappa,\D)$ is a sample of~\eqref{eq:joint} for
geometry $G$, then $(\y,R\bm\kappa,R\D)$ is an exact sample for geometry $RG$. Each training event is presented with a
fresh Haar-random $R$ every time it is drawn (the $Q$ factor of a Gaussian matrix with signs fixed and $\det R=+1$),
so the network is trained on the mixture of all detector orientations, of which the test geometry ($R=I$) is one
member. Distances, times, hit counts, the energy and $\cos\psi_d$ are all invariant under the joint rotation of
positions and axis, so of the inputs only $\bm\xi_d$ changes; in the code the axis therefore need not be rotated at
all, only the positions are. Without the augmentation the validation NLL reached its minimum after a few epochs
and rose afterwards while the training NLL kept falling (memorisation); with it the two stay together.

\section{Outputs, derived quantities and validation}

\paragraph{Back-transformation.}
With $\bm\mu_{\bm z},\Sigma_{\bm z}$ the head's outputs in the standardised variables,
\begin{equation}
  \bm\mu_{\y} = \bm m + \diag(\bm s)\,\bm\mu_{\bm z},\qquad \Sigma_{\y} = \diag(\bm s)\,\Sigma_{\bm z}\,\diag(\bm s) .
\end{equation}
Linear functionals of a Gaussian are Gaussian: for the mean shower depth $\ln(ab) = y_1+y_2$,
\begin{equation}
  \E[\ln ab\mid\D,\bm\kappa] = \mu_1+\mu_2,\qquad \Var[\ln ab\mid\D,\bm\kappa] = \Sigma_{11}+\Sigma_{22}+2\Sigma_{12},
\end{equation}
and $a=e^{y_1}$ is log-normal with median $e^{\mu_1}$ and mean $e^{\mu_1+\Sigma_{11}/2}$. The correlation
$\rho=\Sigma_{12}/\sqrt{\Sigma_{11}\Sigma_{22}}$ measures the $a$--$b$ degeneracy the light leaves.

\paragraph{Resolution.}
The median of $|y_j-\mu_j|$ per component, compared with the median of $|y_j-\bar y_j|$ obtained by predicting
the training mean (the prior). By the law of total covariance,
\begin{equation}
  \Cov(\bm Y) = \E\big[\Cov(\bm Y\mid\D,\bm\kappa)\big] + \Cov\big(\E[\bm Y\mid\D,\bm\kappa]\big),
\end{equation}
so every estimator $\hat{\y}(\D,\bm\kappa)$ satisfies
$\E\|\bm Y-\hat{\y}\|^2\ge\E\big[\tr\Cov(\bm Y\mid\D,\bm\kappa)\big]$: fluctuations that leave no trace in the light,
and the label noise of Remark~1, set the floor.

\paragraph{Calibration.}
If $\bm Y\mid\D,\bm\kappa\sim\N(\bm\mu,\Sigma)$, then for held-out events
\begin{equation}
  r_e^2 = (\y_e-\bm\mu_e)^\top\Sigma_e^{-1}(\y_e-\bm\mu_e)\ \sim\ \chi^2_2,\qquad
  \Pr\big(r_e^2\le\chi^2_{2,\alpha}\big) = \alpha\quad\forall\alpha\in(0,1),
\end{equation}
so 68.3\,\% of the events should fall below $\chi^2_{2,0.683}=2.28$ and 95\,\% below $\chi^2_{2,0.95}=5.99$; the
per-component pulls $(y_j-\mu_j)/\sqrt{\Sigma_{jj}}$ should be standard normal; and the error should grow with the
reported $\sqrt{\Sigma_{jj}}$.

\section{Results}

\paragraph{Stage 1.}
Run of 2026-09-29 on the combined sample (8\,238 training, 2\,059 validation events), vertex, energy and axis
given, 150 epochs requested and stopped at 91; the best epoch (50) has a validation NLL of $1.41$ in the
standardised variables against $\ln2\pi+1=2.84$ for predicting the prior.

\begin{center}
\begin{tabular}{lccc}
\toprule
 & $\ln a$ & $\ln b$ & $\ln(ab)$\\
\midrule
median $|y-\mu|$, posterior mean   & 0.165 & 0.202 & 0.109\\
median $|y-\mu|$, predict the prior mean & 0.277 & 0.210 & 0.269\\
mean $|y-\mu|$, posterior mean     & 0.223 & 0.268 & 0.160\\
median predicted $\sqrt{\Sigma_{jj}}$ & 0.253 & 0.299 & 0.164\\
\bottomrule
\end{tabular}
\end{center}
Coverage $0.73$ / $0.95$ at the $68.3\,\%$ / $95\,\%$ levels: the posterior is calibrated, slightly
conservative. Knowing the vertex, the energy and the axis, the light determines the mean depth $ab$ 2.5 times
better than the prior and $a$ 1.7 times better; $\ln b$ alone is barely improved over the prior, and the
predicted correlation between $\ln a$ and $\ln b$ is strongly negative (median $\rho\approx-0.8$): the
$a$--$b$ degeneracy. Compared with the same network on NC events alone with neither energy nor axis given
(median errors $0.23$ / $0.29$ / $0.18$), the known quantities and the larger sample sharpen every component.

\paragraph{Stage 2, NC/CC from the profile.}
The posterior means were passed to a two-input classifier, $x=(\ln a,\ln b)$, once with the true labels and once
with the inferred values, on the same energy-matched events (1\,234 events, 617 per class, 496 in validation;
matching is needed because $a$ grows with $\ln E$ and the classes differ six-fold in median energy). Across
repeated trainings the true profile gives an AUC of $0.70$--$0.71$ (quadratic discriminant analysis
$0.70$--$0.73$), rising from $0.62$ at $1$--$3$\,GeV to $0.75$ above $10$\,GeV; the inferred profile gives
$0.62$--$0.67$, with the closed-form statistic of the note,
$\Lambda=\ln\N(\hat{\y};\bm m_{\mathrm{CC}},S_{\mathrm{CC}}+\hat\Sigma)-\ln\N(\hat{\y};\bm m_{\mathrm{NC}},S_{\mathrm{NC}}+\hat\Sigma)$,
within $0.01$ of the network in every run. The class difference sits in the shape: NC showers have smaller $a$ and
larger $b$ at the same energy, whereas the mean depth $ab$ alone separates the classes only at the level of
$0.55$--$0.59$. The inferred values move the NC mean toward CC ($\ln\hat b$: $-0.18$ vs $-0.27$ for NC and CC, where
the truth has $-0.10$ vs $-0.30$), which is the mixture-prior shrinkage of Remark~2 acting on the direction that
carries the class information, and it costs the $0.05$--$0.08$ of AUC between the two rows.

\end{document}
